\cotefile{78}
\scanpage{17}
Les données deviennent (en plus de l'entier $n\geq3$)
\[
\left\{\begin{aligned}
 a')\;&\text{Représentation \deleted{fixée}}\\
 &D_n\xrightarrow{\varphi}\mathrm{Aff}(E)\\
 b')\;&s_0\in E
\end{aligned}\right.\tag*{(11)}
\]
\margin{$\begin{aligned}
(D_n&=\underset{\substack{\rotatebox{90}{$\in$}\\\sigma}}{\{\pm1\}}\cdot
\underset{\substack{\rotatebox{90}{$\in$}\\u}}{\mathbb Z/n\mathbb Z}\\
&=\{\sigma,u\mid\sigma^2=1,u^n=1,\\[-2pt]
&\hspace{24pt}\sigma u\sigma^{-1}=u^{-1}\}\\
&=\{\sigma,\sigma'\}\,\sigma^{2}={\sigma'}^{2}=1\}
\end{aligned}$}

Satisfaisant les conditions (pour un pol. rég. affine \uline{unimodulaire})
\[
\deleted{\unsure{(12)}}\quad
\left\{\begin{aligned}
 &\deleted{u s_0\ne s_0}\\
 &\deleted{\unsure{a)}}\ {\scriptstyle 1^\circ}\quad\sigma s_0=s_0\\
 &\deleted{\unsure{b)}}\ {\scriptstyle 2^\circ}\quad s_0,s_1=u s_0,s_2=u^2s_0\ (=u s_1)\quad\text{non alignés}\ (\Rightarrow u s_0\ne s_0)
\end{aligned}\right.\tag*{(13)}
\]
\[
\deleted{\unsure{(14)}}\ (12)\quad
\left[\begin{aligned}
 &\deleted{\unsure{a)}}\ (\text{Pour mémoire})\quad u\text{ unimodulaire}\\
 &\deleted{\unsure{b)}}\quad \varphi\text{ fidèle}
\end{aligned}\right]\quad\begin{aligned}(2^\circ)\\(1^\circ)\end{aligned}
\]
% slanted
\margin{On dira que $\varphi$ est admissible \deleted{s'il satisfait (12) et} s'il existe $s_0$ satisfaisant les conditions (13) $(2^\circ)$ \deleted{\unsure{1)}}}

\uline{Proposition 2} Pour un plan affine $E$ \deleted{fixé} et un entier $n\geq3$ fixés, la donnée dans $E$ d'un $n$-polygône régulier \deleted{\unsure{P}} épinglé équivaut à la donnée d'un couple $(\varphi\unsure{,|;}s_0)$, où $\varphi:D_n\mathbin{\unsure{-|\to}}\mathrm{Aff}(E)$ est un hom de groupes \deleted{injectif}, et où $s_0\in E$ est tel que $\sigma s_0=s_0$ (i.e. $s_0\in E^\sigma$) \uline{et} $s_0,us_0,u^2s_0$ non alignés \deleted{avec en plus} \deleted{\unsure{la|par}} les conditions (12) \deleted{\unsure{à|:}} $\varphi$ \deleted{injectif} \uline{fidèle} et \uline{$u$} \uline{uni}modulaire.

% margin note written sideways
\margin{\deleted{\uline{NB} Il faut \unsure{pour que} \unsure{conjuguer|comprendre} montrer que ces conditions impliquent \unsure{que} \illegible{1} \unsure{si|i.e.} $(+)$}}

Techniquement, la donnée essentielle est celle de $\varphi$. On verra que \deleted{\unsure{si}} \inserted{\text{pour}} les $\varphi$ \deleted{\unsure{p}} correspondant à des polygônes réguliers, \deleted{\unsure{\(\Rightarrow\)}} $\sigma$ est une \uline{réflexion} i.e. $E^{\sigma\deleted{\unsure{0}}}$ $(=\{x\in E\mid \sigma x=x\}$ est une \uline{droite}
% archive sheet number: 16
